\chapter[Conclusions \& Future Scope]{Conclusions \& Future Scope}{Conclusions \& Future Scope}\label{res}

\noindent
\rule{0.49\textwidth}{0.75pt} $_{\bigcirc}$ \rule{0.49\textwidth}{0.75pt}\\
\section{Conclusions}
TrustChain offers a viable solution to safe identity verification by avoiding the process of exchanging uncoded personal data. The system allows users to be in total control of their identity and the institutions to do verification using proof based mechanisms. This minimizes the risks involved in exposing and abusing data.

Decentralized technologies like blockchain and IPFS provide the necessary transparency, immutability and reliability in identity management. Zero Knowledge Proofs help to enhance privacy by authenticating values without providing sensitive information. This is in line with the current security and data protection needs.

The system also initiates user controlled access management where the permissions could be granted or denied any time. Guardian based recovery can further increase reliability and guardian based recovery can guarantee continuity in case of key loss. This renders the platform to be secure and user friendly.

TrustChain, in general, exhibits a scalable and effective identity verification solution in the real world. It effectively changes the paradigm of sharing data to proving based validation, suitable in its adoption in various fields that need identity management.

It lays a robust basis of trust in online communications without invading users privacy. The strategy is in line with the new global requirements of decentralized identity systems. It also shows how cryptographic techniques can be applied in resolving real world security issues. It makes TrustChain a visionary solution in the changing world of digital identity management.
\noindent
\begin{center}
\rule{0.45\textwidth}{0.75pt} \; $\bigcirc$ \; \rule{0.45\textwidth}{0.75pt}
\end{center}

%\chapter{Conclusions \& Future Scope}



\section{Future Scope}
TrustChain may be expanded by means of connecting with official government databases like UIDAI and ITD to enhance the validity of document authentications. This would facilitate the real time validation of credentials and will increase trust into the system in the institutional use cases. 

In addition, enhanced recovery features like managed wallets and encrypted backups may be included. These upgrades will make them easier to use and less reliant on manual recovery processes and also enhance security.

Other technologies that can be integrated into the system include computer vision to verify documents automatically and extract data. This would assist in authenticating the authenticity of uploaded credentials and lessen reliance on third-party verification measures.

Moreover, large scale deployment can be investigated with scalability and performance optimizations. Increasing interoperability with the current digital identity platforms, as well as raising the user experience, will also contribute to the adoption of TrustChain in practical use.


